% !TeX root = ../../Book4.tex
% 纲要标题：## 第6章　复形、同调与 Euler–Poincaré 不变量
% 纲要标签：b4:ch:05
% 纲要位置：计划纲要.md，第 3216 行。

\chapter{复形、同调与 Euler–Poincaré 不变量}
\label{b4:ch:05}
\BookChapterNumber{b4:ch:05}

\begin{chapterabstract}
本章建立复形、Hom 复形、张量复形与同调，说明长正合列和连接同态的构造，并以 Grothendieck 群表述有界复形的 Euler 特征。
\end{chapterabstract}

\begin{learninggoals}
\begin{itemize}
\item 从给定微分计算圈、边界及同调，并正确转换链指标、上链指标和移位。
\item 区分复形映射与任意分次映射，核对 Hom 与张量复形的微分平方。
\item 构造并计算短正合列的连接同态，解释长正合列的自然性。
\item 用短正合关系构造 \(K_0\)，说明 Euler–Poincaré 等式中的有限性条件。
\end{itemize}
\end{learninggoals}

考虑整数上的两步计算：先乘 \(2\)，再取模 \(2\) 的余类。第二步会把第一步的结果全部送到零，但它还会把哪些元素送到零？答案正好是偶数，故这条序列在中间位置没有留下缺口。若只保留第一步而把末项改成零，中间位置的核变成整个 \(\mathbb Z\)，原来的像仍只有 \(2\mathbb Z\)，缺口便是 \(\mathbb Z/2\mathbb Z\)。复形把连续运算的复合为零作为起点；同调则逐处测量「成为零的元素」与「由前一步产生的元素」之间的差别。这个简单计算也提示，比较两个复形不能只比较它们的各项。

\ProseParagraphBreak
本章需要\chapref{b4:ch:04}的交换范畴、核与余核及蛇引理。模的算例另用第二卷第6.1节的张量积泛性质；有限长度版本用第二卷第12.2节已经证明的长度可加性，有限维向量空间则已足够呈现基本计算。复形的基本运算与长正合列可参阅 Weibel\cite[Sections 1.1--1.3, Theorem 1.3.1, Addendum 1.3.3, Proposition 1.3.4]{Weibel1994HomologicalAlgebra}；\(K_0\) 的关系与 Euler 特征见\cite[Chapter II, Definition 6.1.1, Proposition 6.6]{Weibel2013KBook}。移位和映射锥在文献中有不同符号，本书始终采用 \(C[1]^n=C^{n+1}\)、\(\mathrm{d}_{C[1]}=-\mathrm{d}_C\)。

% !TeX root = ../../Book4.tex
% 纲要标题：### 6.1 复形
% 纲要标签：b4:sec:0501
% 纲要位置：计划纲要.md，第 3218 行。

\section{复形}
\label{b4:sec:0501}

相邻映射的复合为零，只保证像包含在核中。同调把两者的差额组织成逐次数的不变量，链与上链的编号则须先固定转换规则。

% 纲要内容（仅作写作计划，尚非教材正文）：
%
% * 链复形与上链复形
% * 圈（cycle）、边界（boundary）及同调
% * 指标转换与次数平移约定
%

\subsection{链复形与上链复形}
\label{b4:subsec:0501:01}

一个序列若处处正合，核与像就完全一致；但许多计算只能保证相邻两步的复合为零。先保留这一较弱的条件，才有可能记录正合性在哪个次数失败，以及核中还剩下什么。

\begin{definition}\label{b4:def:cochain-complex}
设 \(\mathcal A\) 是交换范畴。给定对象 \(C^n\) 及态射 \(\mathrm{d}_C^n:C^n\to C^{n+1}\)（\(n\in\mathbb Z\)），若每个 \(n\) 都满足 \(\mathrm{d}_C^{n+1}\mathrm{d}_C^n=0\)，则称 \(C=(C^n,\mathrm{d}_C^n)\) 为 \(\mathcal A\) 中的\term[shanglian-fuxing]{上链复形}[cochain complex]，\(\mathrm{d}_C^n\) 称为微分。类似地，对象 \(C_n\) 及态射 \(\partial_n:C_n\to C_{n-1}\) 满足 \(\partial_{n-1}\partial_n=0\) 时，组成\term[lian-fuxing]{链复形}[chain complex]。
\end{definition}
\symidx{\(\mathrm{d}_C^n,\ \partial_n\)：上链复形与链复形的微分}

复形条件在每个次数只保证
\[
\mathrm{d}_C^n\mathrm{d}_C^{n-1}=0
\quad\Longleftrightarrow\quad
\operatorname{im}\mathrm{d}_C^{n-1}\subseteq\ker\mathrm{d}_C^n;
\]
正合还要求这里的像与核是同一个子对象。同调将记录这个包含还差多少。链与上链的差别在指标方向：链微分降一次，上链微分升一次。它们的核、像和正合性问题相同。为统一后面的 Hom 复形、移位和映射锥，本编主要用上链指标；讨论消解或 Tor 时也会使用链指标，并明确作转换。

\begin{example}\label{b4:exm:two-term-complex}
在左 \(R\)-模范畴中，任一同态 \(u:M\to N\) 都给出仅在次数 \(0,1\) 非零的复形 \(M\xrightarrow{u}N\)。此时 \(\mathrm{d}^0=u\)，下一微分是 \(\mathrm{d}^1=0_{N,0}:N\to0\)，所以所需条件为
\[
\mathrm{d}^1\mathrm{d}^0=0_{N,0}\circ u=0.
\]
其余相邻复合也均为零。不要求 \(u^2=0\)，因为 \(u\) 的目标是 \(N\)，而源是 \(M\)，一般根本不能把 \(u\) 与自身复合。相反，若每个次数均放同一模 \(M\)，并在每一步都用同一自同态 \(a\)，便确实需要 \(a^2=0\)。例如 \(R=k[\varepsilon]/(\varepsilon^2)\) 上反复乘 \(\varepsilon\) 给出复形。
\end{example}

\begin{definition}\label{b4:def:bounded-complex}
设 \(C\) 是交换范畴中的上链复形。若存在整数 \(a\)，使 \(n<a\) 时 \(C^n=0\)，则称 \(C\) 下有界；若存在 \(b\)，使 \(n>b\) 时 \(C^n=0\)，则称其上有界；同时满足二者时称其\term[youjie-fuxing]{有界复形}[bounded complex]。对象 \(M\) 放在次数零、其余项为零所得复形记为 \(M[0]\)。
\end{definition}
\symidx{\(M[0]\)：集中在零次的复形}

这里的上下指整数次数的大小，与箭头在版面上朝哪个方向无关。

\subsection{圈、边界与同调}
\label{b4:subsec:0501:02}

\begin{definition}\label{b4:def:cohomology}
设 \(C\) 是交换范畴 \(\mathcal A\) 中的上链复形。记
\[
Z^n(C)\coloneqq\ker \mathrm{d}_C^n,\qquad
B^n(C)\coloneqq\operatorname{im}\mathrm{d}_C^{n-1}.
\]
记像分解为 \(\mathrm{d}_C^{n-1}=i^n\pi^{n-1}\)，其中 \(\pi^{n-1}:C^{n-1}\to B^n(C)\) 满，\(i^n:B^n(C)\to C^n\) 单。由 \(\mathrm{d}_C^ni^n\pi^{n-1}=0\) 及 \(\pi^{n-1}\) 满，得到 \(\mathrm{d}_C^ni^n=0\)。因此，核嵌入 \(z^n:Z^n(C)\to C^n\) 的泛性质唯一给出 \(j^n:B^n(C)\to Z^n(C)\)，使
\[
z^nj^n=i^n.
\]
因为 \(i^n\) 单，\(j^n\) 也单。这就是余边对象到余圈对象的典范嵌入，于是定义
\[
H^n(C)\coloneqq\operatorname{coker}\bigl(j^n:B^n(C)\longrightarrow Z^n(C)\bigr).
\]
它称为次数 \(n\in\mathbb Z\) 的\term[shangtongdiao]{上同调对象}[cohomology object]。当 \(R\) 为含幺环、\(C\) 为幺性左 \(R\) 模复形时，\(Z^n(C)\) 的元素称\term[yuquan]{余圈}[cocycle]，\(B^n(C)\) 的元素称\term[yubian]{余边}[coboundary]，且 \(H^n(C)=Z^n(C)/B^n(C)\)。对 \(\mathcal A\) 中的链复形 \((C_n,\partial_n)\)，相应定义 \(Z_n(C)=\ker\partial_n\)、\(B_n(C)=\operatorname{im}\partial_{n+1}\)，余核 \(H_n(C)\coloneqq\operatorname{coker}\bigl(B_n(C)\longrightarrow Z_n(C)\bigr)\) 称为次数 \(n\) 的\term[tongdiao]{同调对象}[homology object]；对模链复形，\(Z_n(C)\) 与 \(B_n(C)\) 的元素分别称\term[quan]{圈}[cycle]与\term[bianjie]{边界}[boundary]。
\end{definition}
\symidx{\(Z_n(C),B_n(C),H_n(C)\)：链复形的圈、边界与同调对象}
\symidx{\(Z^n(C),B^n(C),H^n(C)\)：余圈、余边与上同调对象}

在模范畴中，两个余圈代表同一上同调类，恰好表示它们之差是余边：能够由前一步微分产生的改变被忽略了。在一般交换范畴中，相同的思想由上述余核表达，不能预先假定对象带有元素。

\begin{definition}\label{b4:def:acyclic-complex}
设 \(C\) 是交换范畴中的复形。若所有次数的同调对象均为零，则称 \(C\) 为\term[lingtiao-fuxing]{零调复形}[acyclic complex]，也称正合复形，即在每一个次数处都正合。对上链复形，这等价于 \(\operatorname{im}\mathrm{d}_C^{n-1}=\ker\mathrm{d}_C^n\) 对每个 \(n\in\mathbb Z\) 成立，其中等号指同一个子对象；链复形作相应的指标转换。
\end{definition}

\begin{example}\label{b4:exm:integer-three-complex}
取交换群复形
\[
C^0=\mathbb Z\xrightarrow{\ \mathrm{d}^0\ }C^1=\mathbb Z^2
\xrightarrow{\ \mathrm{d}^1\ }C^2=\mathbb Z,
\qquad \mathrm{d}^0(a)=(2a,0),\quad \mathrm{d}^1(x,y)=3y.
\]
复合为零，且
\[
\begin{aligned}
H^0(C)&=0,\\
H^1(C)&=\mathbb Z(1,0)/2\mathbb Z(1,0)\cong\mathbb Z/2\mathbb Z,\\
H^2(C)&=\mathbb Z/3\mathbb Z.
\end{aligned}
\]
中间的核由 \((1,0)\) 生成，前一步只得到其偶数倍；末项没有下一步限制，却须除去 \(3\mathbb Z\)。这两个商分别记录不同次数的非正合性。
\end{example}

\subsection{指标转换与次数平移约定}
\label{b4:subsec:0501:03}

把链复形改记为上链复形时，本书取
\[
C^n=C_{-n},\qquad \mathrm{d}_C^n=\partial_{-n},\qquad H^n(C)=H_{-n}(C).
\]
因为 \(C^{n+1}=C_{-(n+1)}=C_{-n-1}\)，原来的 \(\partial_{-n}:C_{-n}\to C_{-n-1}\) 正好成为 \(\mathrm{d}_C^n:C^n\to C^{n+1}\)。这只是重编号，没有额外的负号；负号出现在下面的移位中。

\begin{definition}\label{b4:def:complex-shift}
设 \(C\) 是交换范畴中的上链复形，\(r\in\mathbb Z\)。规定
\[
C[r]^n=C^{n+r},\qquad \mathrm{d}_{C[r]}^n=(-1)^r\cdot\mathrm{d}_C^{n+r},
\]
所得复形称为 \(C\) 的 \(r\) 次\term[yiwei]{移位}[shift]。对复形之间的零次态射，移位只改指标，不增加符号。
\end{definition}
\symidx{\(C[r]\)：复形的移位，\(C[r]^n=C^{n+r}\)}

无论是否乘上这个整体符号，相邻微分的复合都仍为零；采用 \((-1)^r\) 并不是复形条件本身所强制的。核与像不因微分乘 \(-1\) 改变，所以
\[
H^n(C[r])=H^{n+r}(C),\qquad (C[r])[s]=C[r+s].
\]
原来次数为 \(m\) 的对象，在 \(C[r]\) 中出现于次数 \(m-r\)，因为 \(C[r]^{m-r}=C^m\)。特别地，\(C[1]^{-1}=C^0\)、\(C[1]^0=C^1\)，所以 \([1]\) 把原来的各项向较小次数移动一格。例如 \(M[0][1]\) 的唯一非零项在次数 \(-1\)，不是次数 \(1\)。若沿用链记号，同一约定写成 \((C[r])_n=C_{n-r}\)、\(\partial_{C[r]}=(-1)^r\partial_C\)。阅读采用相反移位记号的文献时，须同时转换次数和微分。

\ProseParagraphBreak
移位的负号使复形能够相容地组成映射锥，并使同伦在移位后仍满足同一形式的公式。此处先固定约定，\secref{b4:sec:0602}将通过微分平方和连接同态的计算说明它的作用。

% !TeX root = ../../Book4.tex
% 纲要标题：### 6.2 复形的运算
% 纲要标签：b4:sec:0502
% 纲要位置：计划纲要.md，第 3224 行。

\section{复形的运算}
\label{b4:sec:0502}

复形之间的态射必须与微分相容，Hom 和张量积的微分还必须消去两个方向的混合项。逐次的核与余核可以沿用交换范畴的构造，分次符号却需要另作核验。

% 纲要内容（仅作写作计划，尚非教材正文）：
%
% * 复形映射、核、余核与短正合列
% * Hom 复形及张量复形的初步定义
% * 分次符号规则
%

\subsection{复形映射、核、余核与短正合列}
\label{b4:subsec:0502:01}

\begin{definition}\label{b4:def:complex-map}
设 \(C,D\) 是交换范畴 \(\mathcal A\) 中的上链复形。态射族 \(f^n:C^n\to D^n\) 若满足
\[
\mathrm{d}_D^nf^n=f^{n+1}\mathrm{d}_C^n\qquad(n\in\mathbb Z),
\]
则称为\term[fuxing-yingshe]{复形映射}[map of complexes]，记为 \(f:C\to D\)。以逐次复合作为复合，得到复形范畴 \(\operatorname{Ch}(\mathcal A)\)。
\end{definition}
\symidx{\(\operatorname{Ch}(\mathcal A)\)：上链复形及复形映射组成的范畴}

记余圈嵌入为 \(z_X^n:Z^n(X)\to X^n\)（\(X=C,D\)）。交换等式给出 \(\mathrm{d}_D^nf^nz_C^n=f^{n+1}\mathrm{d}_C^nz_C^n=0\)，所以核的泛性质唯一产生
\[
Z^n(f):Z^n(C)\longrightarrow Z^n(D),
\qquad z_D^nZ^n(f)=f^nz_C^n.
\]
另一方面，\(f^n\mathrm{d}_C^{n-1}=\mathrm{d}_D^{n-1}f^{n-1}\) 保证 \(f^n\) 将前一微分的像送入目标的余边对象，因而 \(Z^n(f)\) 将 \(B^n(C)\) 的嵌入送入 \(B^n(D)\)。记 \(q_X^n:Z^n(X)\to H^n(X)\) 为余核态射，则 \(q_D^nZ^n(f)\) 杀死 \(B^n(C)\) 的嵌入，沿 \(q_C^n\) 唯一下降为
\[
H^n(f):H^n(C)\longrightarrow H^n(D),
\qquad H^n(f)q_C^n=q_D^nZ^n(f).
\]
在模范畴中，这两次因子化正是先将余圈 \(z\) 送到 \(f^n(z)\)，再取类 \([z]\mapsto[f^n(z)]\)。恒等映射和复合的保持也由这两种唯一性得到，因此 \(H^n\) 是函子。

\begin{proposition}\label{b4:prop:complex-category-abelian}
若 \(\mathcal A\) 是交换范畴，则 \(\operatorname{Ch}(\mathcal A)\) 也是交换范畴。核、余核、像与有限双积均逐次数计算；特别地，复形序列 \(0\to A\to B\to C\to0\) 短正合，当且仅当每个次数的对象序列短正合。
\end{proposition}
\begin{proof}
给定复形 \(C,D\)，有
\[
\operatorname{Hom}_{\operatorname{Ch}(\mathcal A)}(C,D)
\subseteq\prod_{n\in\mathbb Z}\operatorname{Hom}_{\mathcal A}(C^n,D^n).
\]
由于 \(\mathcal A\) 局部小，右端是集合指标的集合直积，故左端也是集合。这证明了 \(\operatorname{Ch}(\mathcal A)\) 的局部小性。态射逐次相加，零复形和逐次双积给出加性结构。

对 \(f:C\to D\)，令 \(K^n=\ker f^n\)，核嵌入为 \(k^n:K^n\to C^n\)。微分 \(\mathrm{d}_C^n\) 接在这个嵌入之后，得到从 \(K^n\) 到 \(C^{n+1}\) 的态射，且
\[
f^{n+1}\mathrm{d}_C^nk^n=\mathrm{d}_D^nf^nk^n=0.
\]
因此这条态射唯一地通过 \(f^{n+1}\) 的核分解，产生 \(\mathrm{d}_K^n:K^n\to K^{n+1}\)，满足 \(k^{n+1}\mathrm{d}_K^n=\mathrm{d}_C^nk^n\)。再计算
\[
k^{n+2}\mathrm{d}_K^{n+1}\mathrm{d}_K^n
=\mathrm{d}_C^{n+1}\mathrm{d}_C^nk^n=0,
\]
由 \(k^{n+2}\) 单，得到 \(\mathrm{d}_K^{n+1}\mathrm{d}_K^n=0\)。若复形映射 \(h:T\to C\) 满足 \(fh=0\)，其逐次核分解在复合 \(k^{n+1}\) 后由 \(h\) 的微分相容性互相匹配；消去这个单态射，便知分解本身是复形映射。故所得 \(K\to C\) 满足复形范畴中的核泛性质。

对余核 \(q^n:D^n\to Q^n\)，有 \(q^{n+1}\mathrm{d}_D^nf^n=q^{n+1}f^{n+1}\mathrm{d}_C^n=0\)。所以 \(q^{n+1}\mathrm{d}_D^n:D^n\to Q^{n+1}\) 杀死 \(f^n\)，沿其余核 \(q^n\) 唯一下降，产生 \(\mathrm{d}_Q^n:Q^n\to Q^{n+1}\)，满足
\[
\mathrm{d}_Q^nq^n=q^{n+1}\mathrm{d}_D^n.
\]
于是 \(\mathrm{d}_Q^{n+1}\mathrm{d}_Q^nq^n=q^{n+2}\mathrm{d}_D^{n+1}\mathrm{d}_D^n=0\)，消去满态射 \(q^n\) 得到复形条件。任何杀死 \(f\) 的复形映射都逐次通过 \(q^n\) 唯一分解；所得分解的微分相容性也在复合 \(q^n\) 后成立，再由其满性得到。因此 \(D\to Q\) 是复形范畴中的余核。

像是余核的核，余像是核的余核，所以二者也逐次计算。记 \(P=\operatorname{coim}f\)、\(J=\operatorname{im}f\)，结构复形映射为 \(p:C\to P\)、\(i:J\to D\)。每个次数的典范比较 \(a^n:P^n\to J^n\) 满足 \(i^na^np^n=f^n\)，且因 \(\mathcal A\) 交换而可逆。由 \(p,i,f\) 的微分相容性，得到
\[
\begin{aligned}
i^{n+1}\mathrm{d}_J^na^np^n
&=\mathrm{d}_D^nf^n=f^{n+1}\mathrm{d}_C^n\\
&=i^{n+1}a^{n+1}\mathrm{d}_P^np^n.
\end{aligned}
\]
消去单态射 \(i^{n+1}\) 与满态射 \(p^n\)，可知 \(\mathrm{d}_J^na^n=a^{n+1}\mathrm{d}_P^n\)，故这些比较组成复形映射。等式两侧再复合相应的逆，得到
\[
\mathrm{d}_P^n(a^n)^{-1}=(a^{n+1})^{-1}\mathrm{d}_J^n.
\]
所以逐次的逆也组成复形映射，比较确为复形同构。最后的正合性判据来自逐次核和像的识别。
\end{proof}

后面讨论长正合列与导出函子，还需要同调保持 Hom 群的加法。对复形映射 \(f,g:C\to D\)，记它们在余圈对象上诱导的态射为 \(Z^n(f),Z^n(g)\)。核的唯一分解给出 \(Z^n(f+g)=Z^n(f)+Z^n(g)\)。记余核态射为 \(q_X^n:Z^n(X)\to H^n(X)\)（\(X=C,D\)），则
\[
\begin{aligned}
H^n(f+g)q_C^n
&=q_D^n\bigl(Z^n(f)+Z^n(g)\bigr)\\
&=\bigl(H^n(f)+H^n(g)\bigr)q_C^n.
\end{aligned}
\]
由于 \(q_C^n\) 满，得到 \(H^n(f+g)=H^n(f)+H^n(g)\)。所以每个 \(H^n:\operatorname{Ch}(\mathcal A)\to\mathcal A\) 都是加性函子。

\ProseParagraphBreak
逐次分裂的短正合列不必有复形截面。微分可能把各次数所选补对象送出补对象；\secref{b4:sec:0503}中的非零连接同态能够检测这类不分裂现象。不过，连接同态全部为零一般也不足以推出存在复形截面。

\subsection{Hom 复形与张量复形}
\label{b4:subsec:0502:02}

给两个复形逐项取 Hom 或张量积，并不能只把所有微分原样相加：两个方向的混合项必须抵消。以下分别固定产生这种抵消的符号。

\ProseParagraphBreak
次数增加 \(m\) 的分次同态需要同时给出每个次数上的分量 \(f^p:C^p\to D^{p+m}\)，不要求只有有限多个分量非零，因此这些同态组成直积。例如取每个 \(p\in\mathbb Z\) 处都有 \(C^p=\mathbb Z\)，并令微分为零，则恒等映射 \(1_C=(1_{C^p})_{p\in\mathbb Z}\) 有无穷多个非零分量。它属于 \(\prod_p\operatorname{Hom}_{\mathbb Z}(C^p,C^p)\)，却不属于相应直和。

\ProseParagraphBreak
Hom 微分应把次数 \(m\) 的分次同态送到次数 \(m+1\)，而在次数零，微分为零应当恰好表示 \(\mathrm{d}_Df=f\mathrm{d}_C\)。因此取目标微分后的复合与源微分前的复合之差；为了让两次微分中的混合项抵消，第二项的系数随次数增加而变号，得到下面的带符号公式。

\begin{definition}\label{b4:def:hom-complex}
设 \(R\) 是含幺环，\(C,D\) 是左 \(R\)-模的上链复形。定义交换群复形
\[
\operatorname{Hom}_R^m(C,D)\coloneqq
\prod_{p\in\mathbb Z}\operatorname{Hom}_R(C^p,D^{p+m}),
\qquad
(\mathrm{d}f)^p=\mathrm{d}_D^{p+m}f^p-(-1)^m f^{p+1}\mathrm{d}_C^p.
\]
它称为\term[Hom-fuxing]{Hom 复形}[Hom complex]。这里 \(m\in\mathbb Z\)，\(f=(f^p)_{p\in\mathbb Z}\in\operatorname{Hom}_R^m(C,D)\) 是一整族使次数增加 \(m\) 的同态，不要求其与微分交换。
\end{definition}
\symidx{\(\operatorname{Hom}_R^\bullet(C,D),\ \operatorname{Hom}_R^m(C,D)\)：Hom 复形及其第 \(m\) 次项}

\(\mathrm{d}f\) 的次数已是 \(m+1\)，所以再次施行微分时使用的符号是 \((-1)^{m+1}\)。暂略各分量的指标，展开为
\[
\begin{aligned}
\mathrm{d}(\mathrm{d}f)
={}&\mathrm{d}_D^2f-(-1)^m\cdot\mathrm{d}_Df\mathrm{d}_C\\
&-(-1)^{m+1}\cdot\mathrm{d}_Df\mathrm{d}_C
+(-1)^{2m+1}f\mathrm{d}_C^2=0.
\end{aligned}
\]
次数零的余圈正是复形映射。若 \(\mathcal A\) 为一般局部小交换范畴，也可把上述 Hom 换为其态射交换群，得到同样的交换群复形。

\begin{definition}\label{b4:def:tensor-complex}
设 \(R\) 是含幺环，\(C\) 是右 \(R\)-模的上链复形，\(D\) 是左 \(R\)-模的上链复形。定义交换群复形
\[
\begin{gathered}
(C\otimes_R D)^n\coloneqq\bigoplus_{p+q=n}C^p\otimes_R D^q,\\
\mathrm{d}(x\otimes y)=\mathrm{d}_Cx\otimes y+(-1)^p x\otimes \mathrm{d}_Dy,\\
x\in C^p,\quad y\in D^q,\quad p,q\in\mathbb Z.
\end{gathered}
\]
它称为\term[zhangliang-fuxing]{张量复形}[tensor product of complexes]。若 \(R\) 交换且两者均为 \(R\)-模复形，结果也带自然的 \(R\)-模结构。
\end{definition}
\symidx{\(C\otimes_RD\)：复形的张量积}

每个张量项由整数对 \((p,q)\) 标记，取 \(p+q=n\) 的全部项便是总次数 \(n\)。\(\mathrm{d}_C\) 将 \(p\) 增一，\(\mathrm{d}_D\) 将 \(q\) 增一，所以二者都将总次数增一。虽然同一个总次数中可能有无限多个非零项，直和的定义只允许每个元素使用其中有限多个；微分后仍只有有限多个分量。公式与张量积的平衡关系相容，因为微分均为 \(R\)-线性。

\ProseParagraphBreak
混合项有两种产生次序：先作用 \(\mathrm{d}_C\)，则 \(x\) 的次数已经变为 \(p+1\)，接着作用 \(\mathrm{d}_D\) 时带系数 \((-1)^{p+1}\)；反过来先作用 \(\mathrm{d}_D\)，带系数 \((-1)^p\)，再作用 \(\mathrm{d}_C\)。因此它们相消，完整计算为
\[
\mathrm{d}^2(x\otimes y)=\mathrm{d}_C^2x\otimes y
+\bigl((-1)^{p+1}+(-1)^p\bigr)\cdot\mathrm{d}_Cx\otimes \mathrm{d}_Dy
+x\otimes \mathrm{d}_D^2y=0.
\]
这证明了复形条件。张量与同调能否交换是另一问题：后面的\cref{b4:thm:kunneth-pid}将表明，在主理想整环上即使复形各项自由，同调的张量积仍可能漏掉 Tor 项。

\subsection{分次符号规则}
\label{b4:subsec:0502:03}

在本书所用的分次张量结构中，一个次数为 \(r\) 的齐次因子跨过次数为 \(s\) 的齐次因子时，带系数 \((-1)^{rs}\)；相应齐次映射的张量运算也遵守这一约定。这称为\term[fenci-fuhao-guize]{分次符号规则}[graded sign rule]。它适用于具体的交换或张量操作，不能只因公式中出现两个次数就任意加号；是否与微分相容，仍须从所定义的映射核验。

\begin{proposition}\label{b4:prop:graded-leibniz}
设 \(R\) 为含幺环，\(C,D,E\) 是幺性左 \(R\) 模的上链复形，\(r,s\in\mathbb Z\)，\(f\in\operatorname{Hom}_R^r(C,D)\)、\(g\in\operatorname{Hom}_R^s(D,E)\)。以 \((gf)^p=g^{p+r}f^p\)（\(p\in\mathbb Z\)）定义复合，则
\[
\mathrm{d}(gf)=(\mathrm{d}g)f+(-1)^s g(\mathrm{d}f).
\]
这里 \(gf=g\circ f\)，即先作用 \(f\)，再作用 \(g\)；因子在式中仍依次写作 \(g,f\)。在通常的分次乘法公式 \(\mathrm{d}(ab)=(\mathrm{d}a)b+(-1)^{|a|}a\cdot\mathrm{d}b\) 中取 \(a=g\)、\(b=f\)，左因子的次数便是 \(s\)，所以这里出现 \((-1)^s\)。
若 \(R\) 交换且 \(C,D\) 为 \(R\)-模复形，则
\[
\tau:C\otimes_R D\longrightarrow D\otimes_R C,
\qquad \tau(x\otimes y)=(-1)^{pq}y\otimes x
\quad(x\in C^p,\ y\in D^q)
\]
是复形同构，且 \(\tau_{D,C}\tau_{C,D}=1\)。
\end{proposition}
\begin{proof}
第一式右端展开为
\[
\begin{aligned}
(\mathrm{d}g)f+(-1)^s g(\mathrm{d}f)
={}&(\mathrm{d}_Eg-(-1)^sg\mathrm{d}_D)f\\
&+(-1)^sg(\mathrm{d}_Df-(-1)^rf\mathrm{d}_C)\\
={}&\mathrm{d}_Egf-(-1)^{r+s}gf\mathrm{d}_C.
\end{aligned}
\]
这正是左端。第二式的两种复合在纯张量上分别为
\[
\begin{aligned}
\tau\mathrm{d}(x\otimes y)
={}&(-1)^{(p+1)q}y\otimes\mathrm{d}_Cx\\
&+(-1)^{p+p(q+1)}\cdot\mathrm{d}_Dy\otimes x,\\
\mathrm{d}\tau(x\otimes y)
={}&(-1)^{pq}\cdot\mathrm{d}_Dy\otimes x\\
&+(-1)^{pq+q}y\otimes\mathrm{d}_Cx.
\end{aligned}
\]
\(y\otimes\mathrm{d}_Cx\) 的两份系数相等，因为 \((p+1)q=pq+q\)；\(\mathrm{d}_Dy\otimes x\) 的两份系数相等，因为 \(p+p(q+1)=pq+2p\)。因此 \(\tau\) 与微分相容。再交换一次乘上 \((-1)^{2pq}=1\)，所以映射可逆。
\end{proof}

例如两个只在次数一非零的 \(k\)-向量空间复形交换时，交换映射带负号。若删掉该号，在存在非零微分的复形上就可能不再与微分交换。特征二时正负号数值相同，但定义仍沿用同一公式，以便不同特征的论证可以比较。

% !TeX root = ../../Book4.tex
% 纲要标题：### 6.3 上同调长正合列
% 纲要标签：b4:sec:0503
% 纲要位置：计划纲要.md，第 3230 行。

\section{上同调长正合列}
\label{b4:sec:0503}

在模复形的短正合列中，右端余圈可以提升到中项，但其原像未必仍是余圈。对这个提升作微分，便得到连接同态及其带来的上同调长正合列。

% 纲要内容（仅作写作计划，尚非教材正文）：
%
% * 连接同态的构造
% * 长正合列及其自然性
% * 显式计算而非仅引用存在性
%

\subsection{连接同态的构造}
\label{b4:subsec:0503:01}

先在模范畴中考察复形的短正合列
\[
0\longrightarrow A\xrightarrow{i}B\xrightarrow{p}C\longrightarrow0.
\]
若 \(c\in C^n\) 是余圈，将它提升到 \(b\in B^n\) 后，\(b\) 未必仍是余圈。然而 \(p(\mathrm{d}b)=\mathrm{d}c=0\)，故 \(\mathrm{d}b\) 来自 \(A^{n+1}\)。连接同态记录的正是这个提升失败的程度。
\[
\begin{tikzcd}[column sep=large,row sep=large]
A^n\arrow[r,hook,"i^n"]\arrow[d,"\mathrm{d}_A^n"']&B^n\arrow[r,two heads,"p^n"]\arrow[d,"\mathrm{d}_B^n"]&C^n\arrow[d,"\mathrm{d}_C^n"]\\
A^{n+1}\arrow[r,hook,"i^{n+1}"']&B^{n+1}\arrow[r,two heads,"p^{n+1}"']&C^{n+1}
\end{tikzcd}
\]

\begin{proposition}\label{b4:prop:connecting-map-construction}
设 \(\mathcal A\) 是交换范畴，\(0\to A\xrightarrow{i}B\xrightarrow{p}C\to0\) 是上链复形的短正合列。每个整数 \(n\) 都有典范态射
\[
\delta^n:H^n(C)\longrightarrow H^{n+1}(A).
\]
在模范畴中，它由下列规则给出：取 \(c\in Z^n(C)\)，选 \(b\in B^n\) 使 \(p(b)=c\)，唯一写成 \(\mathrm{d}_Bb=i(a)\)，则 \(a\in Z^{n+1}(A)\)，并规定 \(\delta^n[c]=[a]\)。
选定原像 \(b\) 后，元素的去向为：
\[
\begin{tikzcd}[column sep=large,row sep=large]
&b\arrow[r,mapsto,"p^n"]\arrow[d,mapsto,"\mathrm{d}_B^n"']&c\arrow[d,mapsto,"\mathrm{d}_C^n"]\\
a\arrow[r,mapsto,"i^{n+1}"']&\mathrm{d}_Bb\arrow[r,mapsto,"p^{n+1}"']&0
\end{tikzcd}
\]
\end{proposition}
\begin{proof}
固定整数 \(n\)。在此证明中将中间复形暂记为 \(E\)，把短正合列写成 \(0\to A\xrightarrow{i}E\xrightarrow{p}C\to0\)。要运用蛇引理构造连接态射，需要一张竖直箭头的核为 \(H^n\)、余核为 \(H^{n+1}\) 的图。为此，对每个上链复形 \(X\) 定义
\[
Q^n(X)\coloneqq\operatorname{coker}\bigl(B^n(X)\hookrightarrow X^n\bigr),
\qquad q_X:X^n\twoheadrightarrow Q^n(X).
\]
在模范畴中，\(Q^n(X)\) 就是 \(X^n/B^n(X)\)。记余圈嵌入为 \(z_X^r:Z^r(X)\hookrightarrow X^r\)。微分的像包含在 \(Z^{n+1}(X)\) 中，并在 \(B^n(X)\) 上为零，因而唯一诱导
\[
\bar{\mathrm{d}}_X:Q^n(X)\longrightarrow Z^{n+1}(X),
\qquad
\mathrm{d}_X^n=z_X^{n+1}\bar{\mathrm{d}}_Xq_X.
\]

先识别这条态射的核。令 \(k_X:K_X\hookrightarrow Q^n(X)\) 为 \(\bar{\mathrm{d}}_X\) 的核。因为 \(\bar{\mathrm{d}}_Xq_Xz_X^n=0\)，存在唯一 \(u_X:Z^n(X)\to K_X\) 使 \(k_Xu_X=q_Xz_X^n\)。由 \(z_X^{n+1}\) 单，对任意 \(x:T\to X^n\)，条件 \(\bar{\mathrm{d}}_Xq_Xx=0\) 等价于 \(\mathrm{d}_X^nx=0\)，也就是 \(x\) 唯一通过 \(z_X^n\) 分解。若还给定 \(y:T\to K_X\) 满足 \(q_Xx=k_Xy\)，则该分解通过 \(u_X\) 所得的态射与 \(y\) 相等，因为 \(k_X\) 单。这说明 \(Z^n(X)\) 是 \(q_X\) 沿 \(k_X\) 的拉回，故由\cref{b4:lem:epi-pullback}，\(u_X\) 满。又因 \(k_X\) 单，\(u_Xv=0\) 等价于 \(q_Xz_X^nv=0\)，后者恰表示 \(z_X^nv\) 通过 \(B^n(X)\hookrightarrow X^n\) 分解。消去 \(z_X^n\) 可知，\(u_X\) 的核就是 \(B^n(X)\hookrightarrow Z^n(X)\)。由\cref{b4:prop:abelian-normal-factorization}，满态射 \(u_X\) 是其核的余核，故由上同调的定义得到典范识别
\[
\ker\bar{\mathrm{d}}_X\cong H^n(X).
\]
另一方面，\(q_X\) 满，故 \(\bar{\mathrm{d}}_X\) 的像与微分在 \(Z^{n+1}(X)\) 中的像相同，即 \(B^{n+1}(X)\)。因此
\[
\operatorname{coker}\bar{\mathrm{d}}_X
\cong\operatorname{coker}\bigl(B^{n+1}(X)\hookrightarrow Z^{n+1}(X)\bigr)
=H^{n+1}(X).
\]

短正合列中的复形映射保持余圈与余边，因而诱导交换图
\[
\begin{tikzcd}[column sep=large]
Q^n(A)\arrow[r,"Q^n(i)"]\arrow[d,"\bar{\mathrm{d}}_A"']&Q^n(E)\arrow[r,"Q^n(p)"]\arrow[d,"\bar{\mathrm{d}}_E"]&Q^n(C)\arrow[r]\arrow[d,"\bar{\mathrm{d}}_C"]&0\\
Z^{n+1}(A)\arrow[r,"Z^{n+1}(i)"']&Z^{n+1}(E)\arrow[r,"Z^{n+1}(p)"']&Z^{n+1}(C)&
\end{tikzcd}
\]
下排左端单，因为 \(i^{n+1}\) 与余圈嵌入均单。为验证中项正合，设 \(e:T\to Z^{n+1}(E)\) 满足 \(Z^{n+1}(p)e=0\)。由 \(i^{n+1}=\ker p^{n+1}\)，存在唯一 \(a:T\to A^{n+1}\) 使 \(i^{n+1}a=z_E^{n+1}e\)。复形映射的相容性给出
\[
i^{n+2}\mathrm{d}_A^{n+1}a
=\mathrm{d}_E^{n+1}z_E^{n+1}e=0.
\]
由 \(i^{n+2}\) 单，\(a\) 是余圈，故它唯一通过 \(Z^{n+1}(A)\) 分解，得到 \(e\) 在下排左端的原像。反向由 \(pi=0\) 给出。

上排右端满：复合 \(q_Cp^n:E^n\to Q^n(C)\) 满，且它在 \(B^n(E)\) 上为零，故满足
\[
q_Cp^n=Q^n(p)q_E.
\]
原复合满，推出 \(Q^n(p)\) 满。

上排中项的正合性用\cref{b4:lem:generalized-element-calculus}检验。以下每次提升都继续预合成所得满态射，并保留原符号；同一等式中的广义元素具有共同定义域。设 \(\xi:T\to Q^n(E)\) 满足 \(Q^n(p)\xi=0\)。由 \(q_E\) 满，在满态射后取 \(e:T'\to E^n\)，使 \(q_Ee=\xi\)。此时 \(q_Cp^ne=0\)，而 \(\ker q_C=B^n(C)=\operatorname{im}\mathrm{d}_C^{n-1}\)，故经进一步的满态射预合成后可取 \(c':T''\to C^{n-1}\)，使 \(p^ne=\mathrm{d}_C^{n-1}c'\)。再由 \(p^{n-1}\) 满，继续预合成并取 \(e'\) 使 \(p^{n-1}e'=c'\)。于是
\[
p^n\bigl(e-\mathrm{d}_E^{n-1}e'\bigr)=0.
\]
由于 \(i^n=\ker p^n\)，存在唯一 \(a\) 使 \(e-\mathrm{d}_E^{n-1}e'=i^na\)。\(q_E\) 消去余边，所以
\[
\xi=q_Ee=q_Ei^na=Q^n(i)q_Aa.
\]
这给出了 \(\xi\) 在共同满态射后的原像；反向仍由 \(pi=0\) 给出，故上排在中项正合。两排满足蛇引理的假设，已识别的竖直核与余核因此给出所需 \(\delta^n\)。

在模中，蛇引理的提升规则正是陈述中的公式。由 \(i(\mathrm{d}a)=\mathrm{d}_E^2e=0\) 和 \(i\) 单得到 \(\mathrm{d}a=0\)。改变提升 \(e\) 为 \(e+i(a_0)\)，\(a\) 改变为 \(a+\mathrm{d}a_0\)，故其上同调类不变。改变余圈代表 \(c\) 为 \(c+\mathrm{d}c_0\) 时，提升 \(c_0\) 为 \(e_0\)，改取 \(e+\mathrm{d}e_0\) 即可，二次微分为零使 \(a\) 不变。加法也与各步相容，得到所声称的同态公式。
\end{proof}

在一般交换范畴中，最后一段仍可作为广义元素的计算方法，只须在需要提升时先预合成满态射。蛇引理已经保证连接态射存在且与选择无关，这种计算说明了它如何作用于代表元。

\subsection{长正合列及其自然性}
\label{b4:subsec:0503:02}

\begin{theorem}[上同调长正合列]\label{b4:thm:cohomology-long-exact}
设 \(\mathcal A\) 是交换范畴，\(0\to A\xrightarrow{i}B\xrightarrow{p}C\to0\) 是 \(\mathcal A\) 中上链复形的短正合列。以 \(\delta^n:H^n(C)\to H^{n+1}(A)\)（\(n\in\mathbb Z\)）表示其典范连接态射。它们组成\term[shangtongdiao-chang-zhenghelie]{上同调长正合列}[long exact cohomology sequence]
\[
\cdots\longrightarrow H^n(A)\xrightarrow{H^n(i)}H^n(B)
\xrightarrow{H^n(p)}H^n(C)\xrightarrow{\delta^n}H^{n+1}(A)
\longrightarrow\cdots.
\]
给定上述复形短正合列到另一条 \(\mathcal A\) 中的复形短正合列 \(0\to A'\to B'\to C'\to0\) 的态射 \((a,b,c):(A,B,C)\to(A',B',C')\)，记后者的连接态射为 \(\delta'^n\)。诱导的上同调态射与长正合列的所有箭头交换，尤其
\[
H^{n+1}(a)\,\delta^n=\delta'^n H^n(c).
\]
把相邻的四项一起排列，得到自然性的交换图：
\[
\begin{tikzcd}[column sep=large,row sep=large]
H^n(A)\arrow[r,"H^n(i)"]\arrow[d,"H^n(a)"']&H^n(B)\arrow[r,"H^n(p)"]\arrow[d,"H^n(b)"]&H^n(C)\arrow[r,"\delta^n"]\arrow[d,"H^n(c)"]&H^{n+1}(A)\arrow[d,"H^{n+1}(a)"]\\
H^n(A')\arrow[r]&H^n(B')\arrow[r]&H^n(C')\arrow[r,"\delta'^n"']&H^{n+1}(A')
\end{tikzcd}
\]
\end{theorem}
\begin{proof}
对\cref{b4:prop:connecting-map-construction}中的图应用蛇引理，得到
\[
H^n(A)\to H^n(B)\to H^n(C)\to
H^{n+1}(A)\to H^{n+1}(B)\to H^{n+1}(C)
\]
在四个内项处正合。这六项中的箭头与原复形所诱导的箭头相同，因为核、像和余核上的态射都是由原来的 \(i,p\) 唯一诱导的。让 \(n\) 遍历整数，相邻六项列的重叠部分相同，因而拼成所述长列，每项都落在某个六项列的内项位置。

短正合列的态射诱导 \(Q^n\) 与 \(Z^{n+1}\) 两排之间的交换图。\cref{b4:prop:snake-naturality}立即给出连接箭头的交换；其余箭头交换由上同调的函子性得到。这也说明连接同态与次数的排列没有额外选择。
\end{proof}

链指标下，同一定理得到\term[tongdiao-chang-zhenghelie]{同调长正合列}[long exact homology sequence]，连接同态为 \(H_n(C)\to H_{n-1}(A)\)。升次与降次的区别只来自 \(H^n=H_{-n}\)，不表示有两种不同的构造。

\begin{corollary}\label{b4:cor:acyclic-two-of-three}
在交换范畴中，复形短正合列 \(0\to A\to B\to C\to0\) 的三个复形若有两个零调，则第三个也零调。若 \(A\) 零调，则 \(H^n(B)\to H^n(C)\) 对每个 \(n\) 均为同构；若 \(C\) 零调，则 \(H^n(A)\to H^n(B)\) 均为同构。
\end{corollary}
\begin{proof}
把相应零项代入长正合列即可。例如 \(H^n(A)=H^{n+1}(A)=0\) 时，相关三项成为 \(0\to H^n(B)\to H^n(C)\to0\)，所以中间箭头既单又满；其他情形使用同一列在相应次数的正合性。
\end{proof}

\subsection{连接同态的显式计算}
\label{b4:subsec:0503:03}

\begin{example}\label{b4:exm:connecting-nonzero}
设 \(k\) 是域，\(B\) 是次数 \(0,1\) 的复形 \(k\xrightarrow{1}k\)，令 \(A\) 只在次数 \(1\) 等于 \(k\)，以恒等映射嵌入 \(B^1\)，再令 \(C=B/A\)。于是 \(C\) 只在次数零等于 \(k\)，每个次数的短正合列都分裂。可是
\[
H^0(C)=k\xrightarrow{\delta^0}H^1(A)=k
\]
是恒等映射：把 \(c\) 提升为 \(B^0\) 中同一个数，其微分仍为 \(c\)。如果复形短正合列有复形截面，便能把余圈提升为余圈，从而 \(\delta^0=0\)，矛盾。逐次分裂因此弱于复形分裂。
\end{example}

\begin{example}\label{b4:exm:connecting-divisibility}
固定整数 \(m\ge2\)，令 \(D\) 为交换群复形 \(\mathbb Z\xrightarrow{\times m}\mathbb Z\)，非零次数为 \(0,1\)。逐次乘 \(m\) 给出短正合列
\[
0\longrightarrow D\xrightarrow{\times m}D\longrightarrow D/mD\longrightarrow0.
\]
模 \(m\) 后微分为零，故 \(H^0(D/mD)=H^1(D/mD)=\mathbb Z/m\)；原复形则有 \(H^0(D)=0\)、\(H^1(D)=\mathbb Z/m\)。将 \(\bar a\in H^0(D/mD)\) 提升为整数 \(a\)，其微分为 \(ma\)，通过左端嵌入乘 \(m\) 取原像得到 \(a\)。因此 \(\delta^0(\bar a)=\bar a\)，长列的非零部分为
\[
0\longrightarrow\mathbb Z/m\xrightarrow{1}\mathbb Z/m
\xrightarrow{0}\mathbb Z/m\xrightarrow{1}\mathbb Z/m\longrightarrow0.
\]
中间的零映射来自乘 \(m\)，并不破坏正合性，因为它的前一个箭头满、后一个箭头单。
\end{example}

更一般地，只要逐次乘 \(m\) 在模复形 \(D\) 上单射，同一短正合列就成立。连接同态要求先取 \(\mathrm{d}b\)，再用已知的整除关系 \(\mathrm{d}b=ma\) 求 \(a\)，最后取其上同调类；直接把 \(\mathrm{d}b\) 模 \(m\) 约化只会得到零，丢掉了需要记录的信息。

% !TeX root = ../../Book4.tex
% 纲要标题：### 6.4 Euler–Poincaré 等式与 Grothendieck 群
% 纲要标签：b4:sec:0504
% 纲要位置：计划纲要.md，第 3236 行。

\section{Euler–Poincaré 等式与 Grothendieck 群}
\label{b4:sec:0504}

边界对象在相邻次数中以相反符号出现，使复形各项的交错和等于同调的交错和。Grothendieck 群把这一计算推广到没有数值维数的交换范畴。

% 纲要内容（仅作写作计划，尚非教材正文）：
%
% * 有限长度或有限维情形的交错和
% * Grothendieck 群 K_0 中的可加不变量
% * 有界复形与其同调的 Euler 特征相等
%
% ---
%

\subsection{有限长度与有限维情形的交错和}
\label{b4:subsec:0504:01}

求出所有同调群有时很费力，但一个交错和却能在求核取商之前计算。原因是每个边界对象在两个相邻次数中各出现一次，且符号相反。

\begin{definition}\label{b4:def:euler-characteristic}
设 \(k\) 是域，\(C\) 是由有限维 \(k\)-向量空间组成的有界复形。规定
\[
\chi(C)\coloneqq\sum_n(-1)^n\dim_k C^n,
\]
称其为\term[Euler-tezheng]{Euler 特征}[Euler characteristic]。设 \(R\) 为含幺环；若 \(C\) 是由有限长度幺性左 \(R\) 模组成的有界上链复形，以 \(\ell_R(C^n)\) 表示第 \(n\) 项的合成长度，则把维数替换为此长度所得的有限交错和称为长度的 Euler 特征。
\end{definition}
\symidx{\(\chi(C)\)：有界有限维复形的 Euler 特征}
\symidx{\(\ell_R(M)\)：模的合成长度}

有界性保证只有有限多个非零项；每项有限维或有限长度，则保证每个被相加的数值都是有限整数。这是两个不同的要求。这里有限长度指具有有限合成列，其长度是合成因子的个数；长度的良定义与短正合列可加性已在第二卷命题12.2.5“长度的可加性”证明。有限维情形只需秩—零化度公式。两种情形都有
\[
0\to U\to V\to W\to0
\quad\Longrightarrow\quad
\lambda(V)=\lambda(U)+\lambda(W),
\]
其中 \(\lambda\) 分别表示维数或长度。因而对有界复形的短正合列，逐次相加即可得到 \(\chi(B)=\chi(A)+\chi(C)\)。

\begin{example}\label{b4:exm:graph-euler}
取一个有限图，顶点数为 \(v\)，边数为 \(e\)，连通分支数为 \(c\)。给每条边定向，在域 \(k\) 上定义链复形 \(C_1=k^e\xrightarrow{\partial}C_0=k^v\)，其余项为零。若用上链指标 \(C^{-n}=C_n\)，则 \((-1)^{-n}=(-1)^n\)，所以其 Euler 特征仍为 \(v-e\)。以 \(\varepsilon_w\) 表示顶点 \(w\) 的基向量，把一条边送到其终点与起点的基向量之差。每条边界在所属连通分支中的坐标和为零，故像包含在各分支的零和子空间中。反过来，在一个分支中选定基准顶点 \(w_0\)，从 \(w_0\) 到任意顶点 \(w\) 的路径按行进方向给各边加上正负号，其边界就是 \(\varepsilon_w-\varepsilon_{w_0}\)。这些差生成该分支的全部零和向量：若 \(\sum_w a_w=0\)，则
\[
\sum_w a_w\varepsilon_w
=\sum_{w\ne w_0}a_w(\varepsilon_w-\varepsilon_{w_0}).
\]
因此两个子空间相等，每个分支贡献的秩是其顶点数减一，总计 \(\operatorname{rank}\partial=v-c\)，从而
\[
\dim_kH_0=c,\qquad \dim_kH_1=e-v+c,
\qquad v-e=\dim_kH_0-\dim_kH_1.
\]
例如三角形的 \(H_1\) 为一维，而加上一条没有构成新回路的悬挂边不改变这个维数。
\end{example}

\subsection{\texorpdfstring{Grothendieck 群 \(K_0\) 中的可加不变量}{Grothendieck 群 K0 中的可加不变量}}
\label{b4:subsec:0504:02}

维数把所有一维向量空间看成同一种贡献；一般交换范畴中可能有几种不能互相替代的简单对象。为了保留这种差别，可以只规定短正合列应满足的加法关系。Euler 交错和还需要取对象类的差，因此不能只在交换幺半群中相加；下面从自由交换群出发，允许有限整数线性组合，包括负系数。

\begin{definition}\label{b4:def:grothendieck-group}
设 \(\mathcal A\) 是本质小的交换范畴，即其对象同构类组成一个集合。取以各同构类 \([M]\) 为基的自由交换群，再除去全部关系
\[
[B]-[A]-[C]
\qquad\bigl(0\to A\to B\to C\to0\;\text{短正合}\bigr)
\]
所生成的子群，所得群记为 \(K_0(\mathcal A)\)，称为 \(\mathcal A\) 的\term[Grothendieck-qun]{Grothendieck 群}[Grothendieck group]。
\end{definition}
\symidx{\(K_0(\mathcal A)\)：交换范畴的 Grothendieck 群}

本质小的条件保证上述自由群以集合为基。对于全部模组成的大范畴，需要先限制大小或选取所关心的本质小子范畴；不能不加说明地对真类取自由群。本节的交换范畴版本对全部短正合列施加关系。环的 \(K_0(R)\) 则以有限生成投射模的同构类为生成元；它们之间态射在模范畴中的核、余核未必仍为有限生成投射模，因而此范畴通常不是交换范畴。三项都是有限生成投射模的短正合列必因商模投射而分裂，所以这里使用的关系是 \([A\oplus C]=[A]+[C]\)。对象范围和所用的正合列均须区分，不能仅凭同一个符号将两种群等同。\Cref{b4:def:B-grothendieck-group}会以指定正合结构统一这两种定义；\cref{b4:thm:B-euler-isomorphism}进一步证明投射模与完美复形的 Euler 类给出相同的 Grothendieck 群。

\begin{proposition}\label{b4:prop:grothendieck-universal}
设 \(\mathcal A\) 是本质小交换范畴，\(G\) 是交换群。对 \(\mathcal A\) 的对象同构类赋值 \(\lambda(M)\in G\)，若每条短正合列 \(0\to A\to B\to C\to0\) 均满足 \(\lambda(B)=\lambda(A)+\lambda(C)\)，则存在唯一群同态 \(\bar\lambda:K_0(\mathcal A)\to G\)，使 \(\bar\lambda([M])=\lambda(M)\)。若 \(\mathcal B\) 也是本质小交换范畴，则每个正合函子 \(F:\mathcal A\to\mathcal B\) 诱导群同态 \(K_0(F):K_0(\mathcal A)\to K_0(\mathcal B)\)，\([M]\mapsto[F(M)]\)。
\end{proposition}
\begin{proof}
由自由交换群的泛性质，生成元上的赋值唯一延拓为群同态。可加性恰好说定义 \(K_0\) 的每个关系都映到零，故唯一地通过商群分解。第二项中正合函子把定义关系送到定义关系，因此同样可以下降。关系 \([0]=[0]+[0]\) 还说明 \([0]=0\)，不必另外添加零对象关系。
\end{proof}

\begin{example}\label{b4:exm:kzero-vectors}
设 \(k\) 为域。有限维 \(k\)-向量空间范畴满足 \([V]=(\dim_kV)[k]\)，因为基给出 \(V\cong k^{\dim V}\)，且直和对应分裂短正合列。维数诱导 \(K_0\to\mathbb Z\)，而 \(1\mapsto[k]\) 是其逆，故 \(K_0\cong\mathbb Z\)。若改为有限维 \(k\times k\)-模，令 \(e_1=(1,0)\)、\(e_2=(0,1)\)，则 \(M=e_1M\oplus e_2M\)。记两种简单模为 \(S_1=k\times0\)、\(S_2=0\times k\)，并令 \(a=\dim_k e_1M\)、\(b=\dim_k e_2M\)。在各分量中选基给出
\[
e_1M\cong S_1^{\oplus a},\qquad
e_2M\cong S_2^{\oplus b},\qquad
[M]=a[S_1]+b[S_2].
\]
模同态保持这两个分量，短正合列逐分量正合，所以两坐标维数诱导 \(K_0\to\mathbb Z^2\)。它与群同态 \((a,b)\mapsto a[S_1]+b[S_2]\) 互为逆：两种简单模分别映到 \((1,0)\)、\((0,1)\)，而上式把任意 \([M]\) 恢复出来。因此 \(K_0\cong\mathbb Z^2\)。总维数只是这两个坐标之和，不能区别两种简单模。
\end{example}

\subsection{有界复形与同调的 Euler 特征}
\label{b4:subsec:0504:03}

\begin{theorem}[Euler–Poincaré 等式]\label{b4:thm:euler-poincare}
设 \(\mathcal A\) 是本质小交换范畴，\(C\) 是其中的有界复形。则在 \(K_0(\mathcal A)\) 中有
\[
\chi_{K_0}(C)\coloneqq\sum_n(-1)^n[C^n]
=\sum_n(-1)^n[H^n(C)].
\]
这称为\term[Euler-Poincare-dengshi]{Euler–Poincaré 等式}[Euler–Poincaré formula]。因此任何取值于交换群的短正合列可加不变量，都使复形与其同调的交错和相等。
\end{theorem}
\symidx{\(\chi_{K_0}(C)\)：取值于 Grothendieck 群的 Euler 类}
\begin{proof}
将微分分解为到像的满态射再接像的包含态射：
\[
\mathrm{d}^n:C^n\xrightarrow{e^n}B^{n+1}(C)
\xrightarrow{i^{n+1}}C^{n+1},
\qquad \mathrm{d}^n=i^{n+1}e^n.
\]
因为 \(i^{n+1}\) 为单态射，\(\ker e^n=\ker\mathrm{d}^n=Z^n(C)\)。因此对每个 \(n\)，有两条短正合列
\[
0\to Z^n(C)\to C^n\xrightarrow{e^n}B^{n+1}(C)\to0,
\qquad
0\to B^n(C)\to Z^n(C)\to H^n(C)\to0.
\]
第二条来自同调的定义。由这两条列得
\[
[C^n]=[B^n(C)]+[H^n(C)]+[B^{n+1}(C)].
\]
取整数 \(a\le b\)，使 \(C^n=0\) 对 \(n<a\) 或 \(n>b\) 成立。于是 \(B^a(C)=\operatorname{im}\mathrm{d}^{a-1}=0\)，且 \(B^{b+1}(C)=\operatorname{im}\mathrm{d}^b=0\)。将上式乘 \((-1)^n\) 并从 \(a\) 到 \(b\) 求和，其中一组边界项重新编号为
\[
\begin{aligned}
\sum_{n=a}^{b}(-1)^n[B^{n+1}(C)]
&=-\sum_{j=a+1}^{b+1}(-1)^j[B^j(C)]\\
&=-\sum_{j=a}^{b}(-1)^j[B^j(C)].
\end{aligned}
\]
最后一个等号用了上述两个端点为零。这组有限和恰与另一组边界项相消，剩下的正是所需公式。最后对该等式应用\cref{b4:prop:grothendieck-universal}的群同态即可。
\end{proof}

\begin{corollary}\label{b4:cor:euler-quasi-invariance}
设 \(k\) 为域、\(R\) 为含幺环。对有限维 \(k\) 向量空间的有界上链复形，或有限长度幺性左 \(R\) 模的有界上链复形，复形各项的维数或长度交错和等于其上同调的相应交错和。若同一类中的复形映射 \(f:C\to D\) 在每个次数诱导上同调同构，则 \(C,D\) 的相应 Euler 特征相等；此类有界零调复形的相应 Euler 特征为零。
\end{corollary}
\begin{proof}
边界、余圈及上同调仍为有限维向量空间或有限长度模，所以刚才的两条短正合列中，各项的维数或长度均有定义。以 \(\lambda\) 表示这一数值，由可加性得
\[
\lambda(C^n)=\lambda(B^n(C))+\lambda(H^n(C))+\lambda(B^{n+1}(C)).
\]
按前一证明的有限重新编号消去边界项，便得到所需等式。上同调同构使等式右端逐项相同；零调时每项为零。
\end{proof}

有限性不能只当作便于书写的条件。例如令 \(C^n=k\)（\(n\ge0\)）、\(C^n=0\)（\(n<0\)），并取 \(\mathrm{d}^{2r}=1_k\)、\(\mathrm{d}^{2r+1}=0\)（\(r\ge0\)），即
\[
0\to k\xrightarrow{1}k\xrightarrow{0}k\xrightarrow{1}k\xrightarrow{0}\cdots,
\]
其中第一个 \(k\) 位于零次。这是零调复形，但 \(1-1+1-1+\cdots\) 不是这里定义的有限交错和。\(K_0(\mathcal A)\) 的群运算只给出有限求和，并未定义无穷级数运算，所以不能将相邻项成对抵消后宣称这个复形的 Euler 类为零。若要给无界复形赋予更广的 Euler 特征，必须另给允许无限求和的结构和条件，或使所求交错和仍只有有限个非零项；本定理并不提供这样的延拓。


\begin{chaptersummary}
同调对象是边界对象嵌入圈对象的余核；对上链复形，所用的是余边对象与余圈对象。在模范畴中，这就是相邻微分的“核除以像”。复形映射在这些对象之间诱导态射，从而使同调具有函子性。复形范畴的正合性逐次判断，但同调函子一般不保持短正合列；缺少的连接箭头由提升余圈、施加微分和取商构造。Hom 复形使用直积，张量复形使用直和；微分中的符号使连续作用产生的混合项相消，分次复合满足带符号的 Leibniz 律，交换底环上的带符号张量交换与微分相容。

\ProseParagraphBreak
有界复形的边界对象在交错和中成对抵消，因此 Euler 特征可以直接由复形各项计算，也可以由同调计算。Grothendieck 群把这种计算表达为普遍的短正合关系，保留了数值维数可能忘掉的信息。下一章将比较复形之间的同伦与拟同构，说明哪些改变保持同调、哪些改变实际上更强。
\end{chaptersummary}

\BookExercises

\begin{exercise}\label{b4:ex:05-periodic}
在每个整数次数放 \(\mathbb Z/12\mathbb Z\)，每个微分均为乘 \(6\)。验证得到复形并计算各阶上同调。把 \(6\) 改为 \(4\) 后还能得到复形吗？
\end{exercise}
\begin{solution}
\(6^2\equiv0\pmod {12}\)，故是复形。乘 \(6\) 的核为 \(\{\bar0,\bar2,\bar4,\bar6,\bar8,\overline{10}\}\)，由 \(\bar2\) 生成，同构于 \(\mathbb Z/6\mathbb Z\)。像为 \(\{\bar0,\bar6\}\)，在此同构下对应 \(\mathbb Z/6\mathbb Z\) 中由 \(\bar3\) 生成的二阶子群。因此每阶上同调为
\[
\ker(\times6)/\operatorname{im}(\times6)
\cong(\mathbb Z/6\mathbb Z)/\langle\bar3\rangle
\cong\mathbb Z/3\mathbb Z.
\]
而 \(4^2\equiv4\not\equiv0\pmod {12}\)，连续两次微分不为零，不能谈其同调。
\end{solution}

\begin{exercise}\label{b4:ex:05-integer-matrices}
计算次数 \(0,1,2\) 的交换群复形
\(\mathbb Z\xrightarrow{a\mapsto(2a,4a)}\mathbb Z^2
\xrightarrow{(x,y)\mapsto6x-3y}\mathbb Z\)
的上同调。指出中间同调的一个生成元。
\end{exercise}
\begin{solution}
复合为 \(12a-12a=0\)。第一个映射单，故 \(H^0=0\)。第二个映射的核是 \(\mathbb Z(1,2)\)，而前一步的像是 \(2\mathbb Z(1,2)\)，故 \(H^1\cong\mathbb Z/2\mathbb Z\)，生成元由 \((1,2)\) 表示。末项上同调是 \(\mathbb Z/3\mathbb Z\)。
\end{solution}

\begin{exercise}\label{b4:ex:05-shift-support}
设 \(M\) 是模，\(M[0]\) 集中在次数零。分别写出 \(M[0][2]\)、\(M[0][-3]\) 的非零次数。若 \(C\) 的非零项仅在 \([-2,4]\)，\(C[3]\) 的非零项可能位于哪里？
\end{exercise}
\begin{solution}
\(M[0][2]\) 只在 \(-2\) 次非零，\(M[0][-3]\) 只在 \(3\) 次非零。\(C[3]^n=C^{n+3}\)，所以可能非零的范围为 \(-5\le n\le1\)。一般移位 \([r]\) 使支集向左移动 \(r\)，并使微分乘 \((-1)^r\)。
\end{solution}

\begin{exercise}\label{b4:ex:05-hom-support}
设 \(M,N\) 是左 \(R\)-模，\(a,b\in\mathbb Z\)。确定 \(\operatorname{Hom}_R^\bullet(M[0][a],N[0][b])\) 的全部项和微分。
\end{exercise}
\begin{solution}
源仅在 \(p=-a\) 非零，目标仅在 \(-b\) 非零，故 \(p+m=-b\) 等价于 \(m=a-b\)。该次数为 \(\operatorname{Hom}_R(M,N)\)，其余项为零，所有微分为零。因而它是 \(\operatorname{Hom}_R(M,N)[0][b-a]\)：移位 \([b-a]\) 把原来的零次项移到 \(-(b-a)=a-b\) 次，与上述计算一致。
\end{solution}

\begin{exercise}\label{b4:ex:05-tensor-integers}
设 \(m,n\) 是非零整数。令 \(P=(\mathbb Z\xrightarrow{m}\mathbb Z)\)、\(Q=(\mathbb Z\xrightarrow{n}\mathbb Z)\) 均放在次数 \(-1,0\)。写出 \(P\otimes Q\) 的微分并计算上同调。
\end{exercise}
\begin{solution}
在次数 \(-1\) 按 \(P^0\otimes Q^{-1}\)、\(P^{-1}\otimes Q^0\) 排列两项。于是复形为
\[
\mathbb Z\xrightarrow{t\mapsto(mt,-nt)}\mathbb Z^2
\xrightarrow{(u,v)\mapsto nu+mv}\mathbb Z.
\]
第一微分的第二分量来自 \(1\otimes\mathrm{d}_Q\)，其符号为 \((-1)^{-1}=-1\)，因为输入的第一因子在 \(P^{-1}\) 中；第二微分的两个输入分别有第一因子次数 \(0\) 和 \(-1\)，但后一输入只有 \(\mathrm{d}_P\otimes1\) 非零，故得到 \(nu+mv\)。令 \(g=\gcd(m,n)>0\)。第一箭头单，第二箭头的像为 \(g\mathbb Z\)，核由 \((m/g,-n/g)\) 生成，前一像为该生成元的 \(g\) 倍。因此 \(H^{-2}=0\)，\(H^{-1}\cong H^0\cong\mathbb Z/g\mathbb Z\)。
\end{solution}

\begin{exercise}\label{b4:ex:05-tensor-wrong-sign}
设 \(k\) 是特征不为二的域，\(C=D=(k\xrightarrow{1}k)\) 均在次数 \(0,1\)。若把张量微分误写成 \(\mathrm{d}(x\otimes y)=\mathrm{d}x\otimes y+x\otimes \mathrm{d}y\)，证明其平方不为零。
\end{exercise}
\begin{solution}
取次数零的基向量 \(e\in C^0\)、\(f\in D^0\)。其二次微分为 \(2\cdot\mathrm{d}_Ce\otimes \mathrm{d}_Df\)，这是次数二的一维空间中的非零向量。正确公式使两条混合路径的系数分别为 \(-1\) 和 \(1\)，从而相消。
\end{solution}

\begin{exercise}\label{b4:ex:05-homology-not-mono}
找出一个逐次单射的交换群复形映射，使其诱导的某个上同调映射不是单射。再找出逐次满射却不在某个上同调次数满射的例子。
\end{exercise}
\begin{solution}
令 \(D=(\mathbb Z\xrightarrow{2}\mathbb Z)\) 在次数 \(0,1\)。逐次乘 \(2\) 是单射，而 \(H^1(D)=\mathbb Z/2\mathbb Z\) 上诱导零映射，不单。逐次商映射 \(D\to D/2D\) 满，但 \(H^0(D)=0\)、\(H^0(D/2D)=\mathbb Z/2\mathbb Z\)，所以零次上同调映射不满。这不违背函子性：函子保持恒等、复合及同构，一般不保持单态射或满态射；\(H^n\) 的加性也不自动保证这种保持性。
\end{solution}

\begin{exercise}\label{b4:ex:05-prescribed-connecting}
给定左 \(R\)-模同态 \(u:M\to N\)，构造逐次分裂的复形短正合列，其连接同态 \(H^0(C)\to H^1(A)\) 恰为 \(u\)。
\end{exercise}
\begin{solution}
令 \(B=(M\xrightarrow{u}N)\) 在次数 \(0,1\)，令 \(A\) 仅在次数一等于 \(N\)，并嵌入 \(B\)，则 \(C=B/A\) 仅在次数零等于 \(M\)。两次数的短列都显然分裂。\(m\in C^0\) 的提升为同一 \(m\in B^0\)，其微分为 \(u(m)\in N\)，所以 \(\delta^0=u\)。因此连接同态从代数可能性上没有额外限制：任意给定模同态都可由这样的复形短正合列实现。
\end{solution}

\begin{exercise}\label{b4:ex:05-modfour-connecting}
取 \(D=(\mathbb Z\xrightarrow{2}\mathbb Z)\) 在次数 \(0,1\)。计算 \(0\to D\xrightarrow{4}D\to D/4D\to0\) 的连接同态，并写出非零部分的长正合列。
\end{exercise}
\begin{solution}
\(H^0(D)=0\)、\(H^1(D)=\mathbb Z/2\mathbb Z\)。在商复形中，\(H^0(D/4D)=\{\bar0,\bar2\}\cong\mathbb Z/2\mathbb Z\)，\(H^1(D/4D)=(\mathbb Z/4\mathbb Z)/2(\mathbb Z/4\mathbb Z)\cong\mathbb Z/2\mathbb Z\)。提升 \(\bar2\) 为整数二，微分为四，除以左端嵌入的四得到一，所以连接同态是上述生成元下的同构。长列为
\[
0\to\mathbb Z/2\mathbb Z\xrightarrow{1}\mathbb Z/2\mathbb Z
\xrightarrow{0}\mathbb Z/2\mathbb Z\xrightarrow{1}\mathbb Z/2\mathbb Z\to0.
\]
\end{solution}

\begin{exercise}\label{b4:ex:05-split-complex-les}
设交换范畴中的复形短正合列 \(0\to A\to B\to C\to0\) 有复形截面 \(s:C\to B\)。证明所有连接同态为零，并说明 \(H^n(B)\) 的分解依赖于什么选择。
\end{exercise}
\begin{solution}
记短正合列的箭头为 \(i:A\to B\)、\(p:B\to C\)。由 \(ps=1_C\) 得 \(H^n(p)H^n(s)=1\)，所以 \(H^n(p)\) 满；长正合列中的 \(\delta^nH^n(p)=0\) 因而给出 \(\delta^n=0\)。截面确定复形同构 \(F_s=(i,s):A\oplus C\to B\)，取上同调便得到 \(H^n(B)\cong H^n(A)\oplus H^n(C)\)。

若截面改为 \(s'=s+it\)，其中 \(t:C\to A\) 是复形映射，则同一对象的两套直和坐标由
\[
F_{s'}^{-1}F_s=
\begin{pmatrix}1_A&-t\\0&1_C\end{pmatrix}
\]
联系，因为 \(F_{s'}\begin{psmallmatrix}1&-t\\0&1\end{psmallmatrix}=(i,s'-it)=F_s\)。在模范畴中，这就是说 \(b=i(a)+s(c)=i(a-t(c))+s'(c)\)。由 \(H^n\) 的加性，上同调中的坐标变换为 \(\begin{psmallmatrix}1&-H^n(t)\\0&1\end{psmallmatrix}\)；模情形即 \((a,c)\mapsto(a-H^n(t)c,c)\)。所以此直和同构依赖所选的截面。
\end{solution}

\begin{exercise}\label{b4:ex:05-euler-ranks}
设域 \(k\) 上的复形仅在次数 \(0,1,2\) 非零，维数依次为 \(3,5,4\)，且两个微分的秩依次为 \(2,2\)。计算各阶同调维数和 Euler 特征。说明给定秩为何与复形条件相容。
\end{exercise}
\begin{solution}
维数分别为 \(3-2=1\)、\((5-2)-2=1\)、\(4-2=2\)。Euler 特征为 \(3-5+4=2=1-1+2\)。第二微分的核维数为 \(5-2=3\)，第一微分要求的像维数为 \(2\)，故可以把这二维像选在三维核中。再选第一微分为到该像的满射，就得到符合给定秩且复合为零的两张线性映射。
\end{solution}

\begin{exercise}\label{b4:ex:05-finite-groups-kzero}
证明有限交换群范畴的 Grothendieck 群是以素数为指标的自由交换群，其中 \([\mathbb Z/p\mathbb Z]\) 对应第 \(p\) 个基向量。
\end{exercise}
\begin{solution}
有限交换群有合成列：对非零群取非零的极小子群 \(H\)。任取 \(0\ne x\in H\)，极小性给出 \(H=\langle x\rangle\)；若 \(|x|=ab\) 且 \(a,b>1\)，则 \(\langle ax\rangle\) 是阶为 \(b\) 的真非零子群，矛盾。因此 \(H\) 的阶为素数，再对商群按群阶归纳。短正合关系因此把任意 \([M]\) 写成 \([\mathbb Z/p\mathbb Z]\) 的有限整数和。另一方面，令 \(v_p(|M|)\) 为群阶中素数 \(p\) 的指数；短正合列满足 \(|B|=|A|\cdot|C|\)，故这些赋值诱导 \(K_0\to\bigoplus_p\mathbb Z\)。它把 \([\mathbb Z/p\mathbb Z]\) 送到对应基向量，与前述生成映射互为逆，因而没有额外关系。
\end{solution}

\begin{exercise}\label{b4:ex:05-countable-swindle}
令 \(\mathcal A\) 为域 \(k\) 上至多可数维向量空间范畴。证明它本质小且为交换范畴，并计算 \(K_0(\mathcal A)\)。为何结果与有限维范畴不同？
\end{exercise}
\begin{solution}
同构类由 \(0,1,2,\ldots,\aleph_0\) 表示，故本质小；核、余核及有限双积仍至多可数维，故是交换范畴。令 \(V\) 有可数无限基，则对任意 \(W\in\mathcal A\)，都有 \(V\oplus W\cong V\)。分裂正合关系给出 \([V]+[W]=[V]\)，从而 \([W]=0\)，故 \(K_0=0\)。这里使用的是一个已经属于该范畴的无限维对象 \(V\) 及一个分裂短正合关系，并未在 \(K_0\) 中作无限求和。有限维对象不能满足 \(V\oplus k\cong V\)，而可数无限维的 \(V\) 能吸收任意至多可数维的 \(W\)，这使后者的类全部为零。整数维数不能延拓为这一范畴上的可加整数值函数，否则 \([k]=0\) 将与维数一矛盾。
\end{solution}

\begin{exercise}\label{b4:ex:05-kzero-more-than-dimension}
设 \(k\) 为域，\(R=k\times k\)，\(S_1=k\times0\)、\(S_2=0\times k\)。令复形 \(C\) 在次数零为 \(S_1\)、次数一为 \(S_2\)，微分为零。比较它的 \(k\)-维数 Euler 特征与 \(K_0\) 中的 Euler 特征。
\end{exercise}
\begin{solution}
数值为 \(1-1=0\)。在\cref{b4:exm:kzero-vectors}的同构 \(K_0\cong\mathbb Z^2\) 下，\(\chi_{K_0}(C)=(1,-1)\ne0\)。总维数对应群同态
\[
\dim_k:\mathbb Z^2\longrightarrow\mathbb Z,
\qquad(a,b)\longmapsto a+b,
\]
而 \((1,-1)\) 是其核中的非零元素。因此数值 Euler 特征把 \(K_0\) 的信息进一步压缩，总维数为零并不表示各种简单对象的贡献分别抵消。
\end{solution}

\begin{exercise}\label{b4:ex:05-euler-shift-tensor}
对有限维 \(k\)-向量空间的有界复形 \(C,D\)，证明 \(\chi(C[r])=(-1)^r\chi(C)\) 及 \(\chi(C\otimes_kD)=\chi(C)\chi(D)\)。这里是否需要先知道张量复形的同调公式？
\end{exercise}
\begin{solution}
第一式令 \(j=n+r\) 重新编号即可。第二式由有限求和得到
\[
\sum_n(-1)^n\sum_{p+q=n}\dim C^p\dim D^q
=\left(\sum_p(-1)^p\dim C^p\right)
 \left(\sum_q(-1)^q\dim D^q\right).
\]
这两式都不需要张量复形的同调公式：计算只用了各项的直和分解、维数对直和的可加性，以及 \(\dim_k(C^p\otimes_kD^q)=\dim_k C^p\,\dim_k D^q\)。Künneth 公式还进一步确定张量复形的同调，是更强的结论。
\end{solution}

\begin{exercise}\label{b4:ex:05-euler-zero-insufficient}
给出有界有限维复形，使 Euler 特征为零但同调不全为零；再说明仅给定 Euler 特征为什么不能恢复各阶同调。
\end{exercise}
\begin{solution}
取次数 \(0,1\) 的 \(k\xrightarrow0k\)。它有 \(H^0=H^1=k\)，而 Euler 特征为零。次数 \(0,1\) 的 \(k\xrightarrow1k\) 则同调全零且有相同 Euler 特征。交错和会把相邻次数的维数抵消，所以不能保留全部逐次信息。
\end{solution}
